\section{Discussion}
\label{sec:discussion}

\textbf{Why does flexibility hurt?} Daily Bitcoin returns have a standard deviation of about 3\%, while any predictable component is, at best, a small fraction of that. With so little signal relative to noise, a model's error is dominated by its variance: every unit of flexibility spent fitting noise in the training folds is paid back out of sample. Models that predict values close to zero (the baselines, heavily shrunk linear models, auto-ARIMA and conservatively regularized boosting such as CatBoost) cannot lose much, while unconstrained learners such as deep networks, fully grown trees and unregularized OLS lose a great deal. This echoes Welch and Goyal's finding for equity premium prediction \cite{welch2008}, where predictors that look useful in sample routinely fail out of sample. It also means that a model that ``wins'' on an in-sample or single-window evaluation is more likely to be the one that overfitted most.

\textbf{Non-stationarity.} Bitcoin's price rose more than tenfold from the start of the evaluation period to its peak. Features measured in price units (moving averages, MACD, ATR, on-balance volume) therefore take values in the test folds that never appeared in training, which tree-based and neural models cannot extrapolate. Section~\ref{sec:robust} shows how much of the damage this explains. It does not, however, create a positive result: even with stationary features, no model beats the random walk.

\textbf{Directional accuracy is a weak and misleading metric.} Accuracy near 50\% is easily reported as ``better than chance'' without a test. Over 2{,}567 forecasts, the 95\% range for a fair coin is roughly 48--52\%, wide enough to contain every model here except the RNN, which falls below it. At longer horizons, overlapping targets inflate the effective significance and a positive drift rewards models biased toward predicting rises, as the 54.4\% Holt--Winters figure illustrates. We recommend always reporting directional accuracy next to the base rate and an $R^2_{OS}$ against the zero forecast.

\textbf{Implications.} For researchers, a credible claim of predictability should (i) evaluate returns, not price levels; (ii) use walk-forward validation with purging; (iii) span several market regimes; (iv) compare against the zero forecast and buy-and-hold; (v) test significance and correct it for every model and setting tried; and (vi) account for transaction costs. For practitioners, the results caution that off-the-shelf models trained on price and volume data are unlikely to beat simply holding the asset.

\textbf{Limitations.} Our conclusions are bounded by our design. We use only OHLCV-derived features. On-chain activity, order-book, derivatives, macro-financial and sentiment data may contain information that price and volume do not, and the codebase already includes fetchers for several of them. We study daily data, a single asset and library-default hyperparameters; nested tuning within each training fold could help some models, although it would also widen the search that our corrections must account for. Statistical models forecast each test block from the end of training rather than updating daily. Deep models were trained on CPU with early stopping and a 50-epoch cap. Finally, our evaluation period is one realized history; a different seven years could favour different models, which is precisely why we report significance rather than rankings.
